\section{Tercera fase: convertir demanda en pedidos, stock y costes}\label{sec:simulacion}

\begin{idea}{Entrada, trabajo y salida}
\textbf{Entran} previsión central, stock inicial, llegadas comprometidas, errores
de previsión y catálogo económico. \textbf{Se simula} mes a mes el movimiento
de cohortes, las ventas, su antigüedad y las decisiones a tres meses.
\textbf{Salen} cantidades físicamente consistentes y una cuenta económica.
Primero se comprueba una trayectoria sin ruido; después se repite el mecanismo.
\end{idea}

\subsection{Pedido, fabricación asignada y llegada: tres fechas}
El caso aproxima el plazo logístico por tres meses. Si $Q_{s,t}$ es el pedido
decidido al cierre del mes $t$ y $L_{s,t+3}$ la llegada correspondiente,
\begin{equation}\label{eq:plazo}
 L_{s,t+3}=Q_{s,t}
\end{equation}
cuando no existe una interrupción. La fabricación se aproxima por la asignación
ejecutada del pedido; no se representa una planta con su propia capacidad y
secuencia diaria.

Al cerrar diciembre de A012, enero y febrero ya están comprometidos. Marzo es
la primera llegada modificable. Una revisión al cerrar enero modifica abril,
no febrero ni marzo. La distinción evita instrucciones como «pedir en septiembre»
cuando en realidad se necesita \emph{recibir} en septiembre.

\begin{figure}[htbp]\centering
\begin{tikzpicture}[x=2.15cm,y=1cm,font=\small,>=Stealth]
\foreach \x/\mes in {0/dic A012,1/ene A013,2/feb A013,3/mar A013,4/abr A013,5/may A013}{
  \node at (\x,0) {\mes};\draw[gris!50] (\x,-.15)--(\x,-.3);}
\draw[gris!50] (0,-.3)--(5,-.3);
\draw[->,azul,thick] (0,-.65)--node[below]{decisión de diciembre: llegada en marzo}(3,-.65);
\draw[->,azul,thick] (1,-1.5)--node[below]{revisión de enero: llegada en abril}(4,-1.5);
\node[draw=gris!40,fill=fondo,align=center,text width=3.4cm] at (1.5,.9)
  {Enero y febrero\\ya vienen en camino};
\end{tikzpicture}
\caption{El plazo de decisión. El stock disponible y los coches en camino no son
la misma cantidad. Los pedidos de octubre--diciembre de A013 llegarán en A014.}
\label{fig:plazo}\fuente{regla de la sección 7 del enunciado y API logística revisada}
\end{figure}

\subsection{Cómo se fija el estado inicial}
El stock al final de A012 es observado: \num{2877} coches. Las llegadas de enero
y febrero no se conocen directamente y se estiman con la media mensual de
\emph{wholesale} de los últimos seis meses, desde que existe actividad en el
segmento. Se redondean a enteros y se suponen iguales para ambos meses.

\tabla{inicial.tex}{Estado al empezar A013. El stock es observado; el tránsito
es estimado. Las coberturas usan la previsión mensual media A013 como denominador.}
{tab:inicial}{\archivo{resultados/simulacion/tablas/estado_inicial_a013.csv}; total calculado con sumas}

El total supuesto en camino para enero--febrero es \num{4220} coches. Por eso la
cobertura del inventario disponible es de 1,8 meses, pero al incorporar tránsito
sube a 4,5. No se puede usar la segunda cifra para afirmar que todos esos coches
pueden venderse hoy. C ilustra el problema contrario: parte ya de un inventario
alto, que limita la necesidad de pedidos aunque siga existiendo demanda.

Las edades iniciales se representan con lotes de 15, 45 y 75 días; 105, 135 y
165 días; y 210 días para el tramo superior a 180. Se conservan unidades enteras.
Las sensibilidades desplazan las edades 15 días y comparan referencias de tránsito
de tres o seis meses, o cantidades un 25\,\% menores o mayores.

\subsection{El balance físico de un mes}
Sin transfers, con demanda realizada $D_{s,t}$ y stock anterior $S_{s,t-1}$:
\begin{align}
 V^{\mathrm{red}}_{s,t}&=\min(D_{s,t},S_{s,t-1}+L_{s,t}),\\
 S_{s,t}&=S_{s,t-1}+L_{s,t}-V^{\mathrm{red}}_{s,t}.
\end{align}
Las llegadas se consideran disponibles al inicio del mes y las ventas se imputan
tras 15 días. Es una convención mensual para valorar edad y coste, no una
simulación de cada comprador durante el día.

\paso{Ejemplo pequeño.}
Hay 20 coches y llegan 50. Si 80 clientes quieren comprar, se venden 70 de la
red y quedan diez sin atender. Con cinco transfers directos, las ventas totales
son 75 y se pierden cinco. El stock final de la red es cero en ambos casos:
los transfers no deben restarse otra vez de un stock al que no se incorporaron.

En la alternativa de transfer se utiliza una fracción $a$ del déficit:
\begin{align}
 T_{s,t}&=\left\lfloor a\,\bigl(D_{s,t}-V^{\mathrm{red}}_{s,t}\bigr)\right\rfloor,\\
 V_{s,t}&=V^{\mathrm{red}}_{s,t}+T_{s,t},\qquad
 P_{s,t}=D_{s,t}-V_{s,t},
\end{align}
con $P$ igual a ventas perdidas. El caso base usa $a=0$ porque no hay
disponibilidad garantizada. Se contrasta $a=0{,}25$ y $a=1$ como referencia
favorable; no se asume que los países europeos tengan siempre stock ilimitado.

\subsection{Por qué FIFO cambia el resultado económico}
El balance total no informa de la edad de los coches que se venden. FIFO
significa que se atiende la demanda con las cohortes más antiguas primero.
Dos políticas con el mismo stock final pueden generar gastos comerciales
distintos si conservan edades diferentes.

En cada mes se añade la cohorte nueva, se venden primero las antiguas, se
imputa el VM de cada edad y se envejece lo que queda. Las cohortes iniciales
mantienen sus días anteriores a enero: un coche con 170 días no recibe otros
180 días subvencionados al cambiar de año.

\subsection{Coste de almacenamiento e intereses}
Para un coche de edad inicial $a$ que permanece otros $d$ días, el coste adicional
de la marca es
\begin{equation}\label{eq:coste-stock}
 C^{\mathrm{marca}}(a,d)
 =2\left[\min(a+d,180)-\min(a,180)\right]\unidad{\euro}.
\end{equation}
Así, un coche recién llegado que se vende en el mismo mes genera
$2\times15=30$ euros. Un coche de 170 días que permanece 30 días adicionales
genera sólo $2\times10=20$ euros. El máximo de toda su vida facturada es
360 euros; no se extrapola el coste durante años.

\begin{figure}[htbp]\centering
\begin{tikzpicture}
\begin{axis}[width=.88\linewidth,height=6cm,xmin=0,xmax=300,ymin=0,ymax=410,
  xlabel={Edad total del coche (días)},ylabel={Coste acumulado de la marca (\euro)},
  xtick={0,60,90,120,180,240,300},ytick={0,120,240,360},legend pos=south east]
\addplot[azul,thick,domain=0:300,samples=101]{2*min(x,180)};
\addlegendentry{Regla del caso}
\addplot[gris,dashed] coordinates {(180,0) (180,360)};
\end{axis}\end{tikzpicture}
\caption{El coste de la marca se detiene a los 180 días. El gasto comercial por
edad y los intereses del concesionario son conceptos adicionales.}
\label{fig:coste-stock}\fuente{enunciado, sección 2.4; ecuación~\eqref{eq:coste-stock}}
\end{figure}

Los intereses del concesionario se calculan por tramos: tarifa diaria del primer
tramo hasta el día 90, del segundo entre 91 y 180, y ningún interés adicional
contado por el caso después. Si $c_1,c_2$ son las tarifas diarias del catálogo,
la función acumulada es
\begin{equation}
 J(a)=c_1\min(a,90)+c_2\min\!\left(\max(a-90,0),90\right).
\end{equation}
El coste de un periodo es $J(a+d)-J(a)$. Por ejemplo, pasar de 85 a 100 días
supone cinco días al primer tipo y diez al segundo, no quince días a ambos.
Estos intereses se muestran separados del resultado de la marca. El vencimiento
del crédito a los 180 días es un riesgo de liquidez para el concesionario.

\subsection{Cuánto pedir y cómo dimensionar el colchón}
Al cerrar $t$, se proyecta el stock disponible antes de la primera llegada
modificable, usando las llegadas ya comprometidas y la previsión de $t+1$ y
$t+2$. La implementación no permite stock proyectado negativo: una demanda que
no se atiende no se convierte en un inventario negativo a reponer, porque el
caso base no acumula pedidos pendientes de clientes.

La regla de pedido es
\begin{equation}\label{eq:pedido}
 Q_{s,t}=\left\lceil\max\!\left(0,
 \widehat D_{s,t+3}+SS_{s,t}-\widehat S_{s,t+2}\right)\right\rceil.
\end{equation}
El redondeo hacia arriba expresa coches enteros. $SS$ es el stock de seguridad;
no es una reserva idéntica para todos los segmentos.

La fórmula docente $z\sigma\sqrt{3}$ supone errores mensuales independientes.
La revisión utiliza errores acumulados observados en ventanas de tres meses:
\begin{equation}
 \sigma_{3,s}^{\mathrm{rel}}
 =\operatorname{sd}\left(
 \frac{\sum_{u=t+1}^{t+3}(\widehat D_{s,u}-O_{s,u})}
      {\sum_{u=t+1}^{t+3}\widehat D_{s,u}}\right),\qquad
 SS_{s,t}=z_s\,\sigma_{3,s}^{\mathrm{rel}}
                \sum_{u=t+1}^{t+3}\widehat D_{s,u}.
\end{equation}
Las ventanas no se unen entre años de validación distintos. El error ya está
acumulado durante el plazo; multiplicar otra vez por $\sqrt3$ duplicaría esa
adaptación. Se cubren los meses \emph{hasta la llegada}, no los tres posteriores
a recibirla.

\paso{Comprobación del ejemplo del enunciado.}
Con 536 coches iniciales, 700 de llegada en enero y otros 700 en febrero,
y previsiones 550 y 600, el stock proyectado es $536+700+700-550-600=786$.
Con $\sigma=100$ y $z=1{,}28$, el colchón docente es aproximadamente 222 coches.
Para una previsión de marzo de 850, resulta
\[
 Q=\left\lceil850+1{,}28\times100\sqrt3-786\right\rceil=286.
\]
El notebook comprueba este valor ejercitando el planificador, además de verificar
balances y fronteras de coste.

\subsection{Revisar sin conocer la demanda que no se atendió}
La revisión utiliza el cociente entre las ventas y la previsión acumuladas en
los tres últimos meses, $r$. La corrección amortiguada es
\begin{equation}\label{eq:revision}
 f=\max\bigl(0,\,1+0{,}5(r-1)\bigr).
\end{equation}
Si las ventas superan un 20\,\% la previsión, el factor es 1,10: se incorpora
la mitad de la desviación para evitar reaccionar plenamente a un mes ruidoso.

Si la red agotó su disponibilidad en la ventana, las ventas están censuradas
por falta de coches. En ese caso el factor no se reduce por debajo de uno.
El planificador observa ventas y disponibilidad, no utiliza mágicamente la
demanda latente simulada para corregirse. El factor se aplica a las previsiones
del plazo y sólo cambia el pedido que llegará en $t+3$.

\subsection{Primero comprobar un año sin ruido}
La referencia determinista utiliza la previsión central redondeada al entero
más próximo, con $z=1{,}28$. No es todavía la política recomendada ni una media
de mil simulaciones.

\subsubsection*{Un recorrido real: SSUV en el primer trimestre}
\begin{table}[H]\centering
\caption{Primer trimestre SSUV de la referencia determinista, coches.}
\begin{tabular}{@{}lrrrrrr@{}}\toprule
Mes & \makecell{Stock\\inicial} & Llegadas & Demanda & Ventas & Pérdidas & \makecell{Stock\\final} \\\midrule
Enero & 0 & 390 & 369 & 369 & 0 & 21 \\
Febrero & 21 & 390 & 500 & 411 & 89 & 0 \\
Marzo & 0 & 1.420 & 591 & 591 & 0 & 829 \\\bottomrule
\end{tabular}
\fuente{\archivo{resultados/simulacion/tablas/simulacion_central_mensual_a013.csv}}
\end{table}

En enero, $0+390-369=21$ coches quedan en stock. En febrero sólo hay
$21+390=411$ para una demanda de 500, de modo que se pierden 89 compras.
La llegada de marzo todavía no está disponible para esos clientes. Al llegar
los \num{1420} coches de marzo, la demanda es 591 y quedan 829 de protección;
las compras perdidas de febrero no se convierten en pedidos pendientes.

El coste de enero también puede comprobarse a mano: 369 coches vendidos con
15 días y 21 que permanecen 31 días generan
$369\times15\times2+21\times31\times2=\num{12372}$ euros. La convención mensual
es la misma que se usa después en todas las réplicas aleatorias.

\tabla{determinista.tex}{Referencia determinista A013: demanda central sin ruido,
colchón uniforme $z=1{,}28$. Resultado tras costes modelizados.}
{tab:determinista}{\archivo{resultados/simulacion/tablas/resumen_central_a013.csv}}

Atiende \num{18991} coches y deja \num{89} sin atender, todos en SSUV.
La contribución es 64,18 M\euro{} y el stock final, \num{2319}. Tener prevista
la demanda no elimina las restricciones iniciales: en SSUV se parte de stock
cero y las dos primeras llegadas ya están comprometidas.

\FloatBarrier
\grafosanchos[\textbf{Antes de añadir azar.} Se verifican balances y unidades
enteras, FIFO y VM por edad, los 15 días de una venta nueva, el límite de 180 días,
intereses por tramos, transfers, ausencia de anticipación y el pedido de diciembre
que llega en marzo A014. Repetir una simulación no corrige un mecanismo mal definido.]
{stock_central_a013_a014.png}
{Referencia determinista y prolongación hasta junio de A014. El punto gris es
stock inicial observado; el trazo azul, stock simulado. La línea vertical separa
A013 de la sensibilidad de continuidad, que no es una previsión validada.}
{fig:simulacion}{\archivo{simulacion.ipynb}; política base, escalas por segmento}
